\documentclass[11pt]{article}
\usepackage[margin=1in]{geometry}
\setlength{\emergencystretch}{3em}
\usepackage{amsmath,amssymb,enumitem}
\usepackage[hidelinks]{hyperref}
\title{Three Ballots in One: Ballot Structure, Coalition Anatomy, and Demographic Splits in the 2026 Alberta Referendum}
\author{Liam McGrath}
\date{25 October 2026}
\begin{document}
\maketitle
\tableofcontents

\section{Overview of the 19 October 2026 Ballot}
\label{sec:ballot-overview}
On 19 October 2026, Alberta voters cast more than $2.07\times 10^6$ ballots across ten provincial referendum questions---the largest ballot count in the province's history. Although presented on a single ballot, the ten questions partitioned into three distinct constitutional and policy regimes:
\begin{enumerate}[label=\arabic*.]
\item \textbf{Questions 1--5 (Non-constitutional provincial policy):} Five proposals on provincial immigration limits, eligibility for provincially funded health and social programs, a 12-month residency waiting period, fees for non-permanent residents, and proof of citizenship to vote.
\item \textbf{Questions 6--9 (Constitutional amendments):} Four proposals requiring bilateral or multilateral amendment of the Constitution of Canada: provincial selection of King's Bench and Appeal judges, abolition of the federal Senate, opting out of federal social programs with full funding, and provincial paramountcy over conflicting federal laws.
\item \textbf{Question 10 (Sovereignty process):} A binary choice between remaining a province of Canada (Option 1: \textit{Remain}) or commencing the legal process toward a binding referendum on separating from Canada (Option 2: \textit{Referendum}).
\end{enumerate}

Let $p_q \in [0,1]$ denote the two-way affirmative share on question $q \in \{1,\dots,10\}$:
\begin{equation}
p_q = \frac{V_q^{\text{For}}}{V_q^{\text{For}} + V_q^{\text{Against}}}.
\label{eq:two-way-share}
\end{equation}
The official Elections Alberta count reveals a clean regime boundary at the $p_q = 0.50$ threshold:
\begin{equation}
\min_{q \in \{1,\dots,5\}} p_q = 0.5229 > 0.50 > \max_{q \in \{6,\dots,9\}} p_q = 0.4978, \qquad p_{10}^{\text{Remain}} = 0.7495.
\label{eq:regime-separation}
\end{equation}

\begin{center}
\begin{tabular}{lp{6.4cm}rrr}
\hline
\textbf{Q\#} & \textbf{Subject} & \textbf{For / Stay} & \textbf{Against} & \textbf{For \%} \\
\hline
Q1 & Provincial immigration control \& job priority & 1,093,541 & 907,731 & 54.64 \\
Q2 & Citizen/PR eligibility for social programs & 1,032,950 & 941,379 & 52.32 \\
Q3 & 12-month residency wait for benefits & 1,053,247 & 932,163 & 53.05 \\
Q4 & Healthcare/school fee for non-perm.\ status & 1,007,376 & 919,120 & 52.29 \\
Q5 & Proof of citizenship required to vote & 1,083,647 & 730,033 & 59.75 \\
\hline
Q6 & Provincial selection of superior-court judges & 832,386 & 918,241 & 47.55 \\
Q7 & Abolish the unelected federal Senate & 898,093 & 926,950 & 49.21 \\
Q8 & Opt out of federal programs w/ full funding & 752,760 & 855,376 & 46.81 \\
Q9 & Provincial law priority over federal law & 785,995 & 792,812 & 49.78 \\
\hline
Q10 & Remain in Canada vs.\ separation referendum & 1,551,868 & 518,766 & 74.95 \\
\hline
\end{tabular}
\end{center}

\section{Ballot Roll-Off and Issue Salience}
\label{sec:roll-off}
Comparing total valid votes $N_q = V_q^{\text{For}} + V_q^{\text{Against}}$ across the ten questions measures voter roll-off relative to Question 10 ($N_{10} = 2{,}070{,}634$):
\begin{equation}
\rho_q = 1 - \frac{N_q}{N_{10}}.
\label{eq:rolloff}
\end{equation}
Roll-off was minimal on Question 1 ($\rho_1 \simeq 3.35\%$) and Question 3 ($\rho_3 \simeq 4.12\%$), rose to $\rho_5 \simeq 12.4\%$ on voter citizenship ID, and peaked on the technical constitutional paramountcy and fiscal opt-out questions ($\rho_8 \simeq 22.3\%$, $\rho_9 \simeq 23.8\%$). Roughly one in five voters who cast a decisive ballot on remaining in Canada declined to endorse constitutional brinkmanship on Questions 6--9.

\section{Demographic Decomposition of Question 10}
\label{sec:demographics}
Combining the Elections Alberta tally with pre-referendum and exit-poll cross-tabs from Abacus Data ($n = 1{,}021$, Oct 7--13, 2026) and Research Co.\ ($n = 811$ likely voters, Oct 16--18; $n = 601$ exit poll, Oct 19--20) shows how the aggregate $74.95\%$ Remain share arises as a turnout-weighted sum across demographic strata $g$:
\begin{equation}
P(\text{Remain}) = \sum_{g} w_g\, P(\text{Remain} \mid g), \qquad \sum_g w_g = 1.
\label{eq:law-total-prob}
\end{equation}

\subsection{Regional Gradient}
Remain won every geographic region of Alberta, but with a 20-point urban--rural gradient:
\begin{itemize}
\item \textbf{Edmonton Metro:} $77\%$ Remain, $17\%$ Option 2, $6\%$ undecided.
\item \textbf{Calgary Metro:} $75\%$ Remain, $19\%$ Option 2, $6\%$ undecided.
\item \textbf{Small Cities \& Suburbs:} $62\%$ Remain, $29\%$ Option 2, $9\%$ undecided.
\item \textbf{Rural Alberta:} $57\%$ Remain, $35\%$ Option 2, $8\%$ undecided.
\end{itemize}

\subsection{Partisan Asymmetry}
The referendum's sharpest cleavage ran inside the governing United Conservative Party (UCP) coalition rather than between the two major parties as wholes:
\begin{itemize}
\item \textbf{Provincial NDP supporters:} $93\%$ Remain, $4\%$ Option 2.
\item \textbf{Provincial UCP supporters:} $44\%$ Remain, $52\%$ Option 2 ($4\%$ undecided), with $77\%$ of all Option 2 voters coming from the UCP base. Among UCP supporters voting Remain, $52\%$ cited economic caution or restraint rather than satisfaction with Ottawa.
\item \textbf{2025 Federal Vote:} Federal Liberal ($95\%$ Remain / $5\%$ Option 2) and Federal NDP ($82\%$ Remain / $18\%$ Option 2) voters broke overwhelmingly for Remain, while 2025 Federal Conservative voters split $65\%$ Remain to $35\%$ Option 2.
\end{itemize}

\subsection{Age, Turnout, and Identity Attachment}
Turnout strongly favored older cohorts: $85\%$ of voters aged $60+$ definitely voted (including mail-in ballots, which broke $93\%$ Remain), compared with $38\%$ of voters under $30$. Older Albertans ($60+$) reported identical $71\%$ strong attachment to both Canada and Alberta. By contrast, interest in Alberta becoming a US state skewed toward younger voters ($21\%$ among ages 18--34 vs.\ $15\%$ among 35--54 and $6\%$ among $55+$).

\section{Decomposing the Option 2 Coalition: Leverage vs.\ Secession}
\label{sec:coalition-math}
A critical analytical mistake is equating the $25.05\%$ Option 2 vote with hard support for independence. Let $M \in \{\text{Independence}, \text{Better Deal}, \text{Direct Say}, \text{Message to Ottawa}, \text{Other}\}$ denote an Option 2 voter's primary motive, with empirical shares $(0.32, 0.23, 0.22, 0.19, 0.04)$, and let $P(\text{Separate} \mid \text{Option 2}) = 0.63$ be the share of Option 2 voters who say they would vote to separate in a binding follow-up referendum. Accounting for non-separatist protest voters yields an implied direct-independence share across the entire electorate of
\begin{equation}
P(\text{Independent Alberta}) \simeq 0.18 \quad (18\%),
\label{eq:direct-independence}
\end{equation}
demonstrating that roughly one-third of the Option 2 vote functioned as constitutional bargaining leverage rather than a mandate to leave Confederation.

\end{document}
